Design Organisational Approval

Privileges: A DOA holder can
Perform design activities within the scope of approval
Have compliance documents accepted by the Agency without further verification
Perform activities independently from the Agency

We are most interested in the first rule, and for appliance we have to read:
nitial Airworthiness Regulation, including Part 21
Relevant section of Regulation: Articles 1(h) & 8
Annex I to Regulation (EC) No 748/2012 - Part 21 - Certification of aircraft and related products, parts and appliances, and of design and production organisations
Relevant section of Part 21: Subpart J - Design Organisation Approval

Relevant link once I find it: https://eur-lex.europa.eu/legal-content/EN/TXT/?uri=CELEX%3A02012R0748-20230825
We need to cover: 
-safety management element
-design assurance element
-safety risk management
-Overview of the entire logistic system, including maintenance, product and parts
- Agency a handbook that describes, directly or by cross reference, the organisation, its relevant policies, processes and procedures (we are delft??)
For the change we just need the following:
After the issue of a design organisation approval, each change to the design management system that is significant to the demonstration of compliance or to the airworthiness, operational suitability and environmental protection of the product, part or appliance shall be approved by the Agency before being implemented. The design organisation shall submit to the Agency an application for approval demonstrating, on the basis of the proposed changes to the handbook, that it will continue to comply with this Annex.

21.A.251 The terms of approval shall identify the types of design work, the categories of products, parts and appliances for which the design organisation holds a design organisation approval, and the functions and duties that the organisation is approved to perform with regard to the airworthiness, operational suitability and environmental characteristics of the products. For design organisation approvals covering type-certification or European Technical Standard Order (ETSO) authorisation for auxiliary power units (APUs), the terms of approval shall contain in addition the list of products or APUs. Those terms shall be issued as part of a design organisation approval.


\begin{comment}
REAT THE DOWN BULLSHIT
\end{comment}
If we apply for this DOA we will get certified so that we will have the capability of internally managing and approving ourselves. This means we as an entity that can approve our own design modification up to a certain extent without the need of an EASA direct approval. We become our own supervisors essentially. 

The benefits that we have after receiving the DOA is that we can internally classify a change as "Minor" or "Major". (KEEP THIS IN MIND) 

The process of documentation and proof of compliance involves the following steps: classifying the action, designing, performing compliance demonstration, validation as well as a statement of compliance, before issuing a STATEMENT OF APPROVAL.

21.A.263
In terms of design changes, we are allowed to perform modifications/repairs that fall under the "MINOR" category ourselves. We just have to follow the steps mentioned above once we get the DOA ourselves.

For modifications/repairs that fall under "Major" class, we can perform all the steps EXCEPT THE LAST ONE, we cannot issue a statement of approval.

Now here is the cool part. It is stated in 21.A.251 Terms of approval, " The terms of approval shall identify the types of design work, the categories of products, parts and appliances for which the design organisation holds a design organisation approval", from my understanding, when creating the DOA, we have to reach an agreement with EASA to consider the wing modifications and other structural modifications as "minor".


Q: Would EASA allow us to classify the wing change as a "Minor" design consideration when we are establishing the DOA?

While reading the documentation for Design Organisational Approval, we found out that if we are approved for this category, we can internally design, validate and certify "minor" changes brought to the design. It is also stipulated that we have the capability to classify ourselves the class of the modification. My question is, when setting up the DOA, can we set it up in such a manner with EASA that the change of the wings or tail would be classified under "minor" for us? Would EASA allow that, and is it possible?
